\section{Discussion}

The central finding is not that academically challenged students seek more help. Instead, different forms of academic challenge are associated with sharply different observable responses. Students who maintain relatively stronger performance while studying in low-pass-rate classrooms are the most likely to use formal mathematics support, whereas students with weak classroom-relative performance---including those simultaneously in difficult classroom contexts---use CMAT much less. The result survives nine pre-specified state definitions and a continuous model, making it unlikely to be a consequence of one convenient cutoff.

This distinction matters for both engagement research and mathematics-support practice. If support-centre use were treated directly as engagement, the contextual-challenge group would appear highly engaged and the individual-strain group less engaged. Kahu's framework cautions against that inference \citep{kahu2013}. We observe one behaviour---formal institutional help-seeking---not the student's cognitive engagement, emotional investment, informal support network, or independent study. The administrative evidence therefore supports a narrower statement: \emph{students experiencing greater academic need are not necessarily the students making greater use of formal support}. That gap is substantively important in its own right.

The finding aligns with prior mathematics-support research on non-engagement. Mac an Bhaird et al. \citep{macanbhaird2013} show that availability of mathematics learning support does not guarantee appropriate uptake, while broader reviews document continuing interest in who engages with support and why others do not \citep{lawson2020,mullen2024}. Our contribution is to add a longitudinal and contextual dimension. Rather than classifying students solely by attendance or prior preparation, we distinguish how they perform relative to classmates from the observed outcome environment of the classroom. The resulting pattern suggests that the students who mobilize formal help successfully may not be those with the weakest realized outcomes.

One possible interpretation is that formal help-seeking requires resources beyond academic need itself: awareness, confidence, familiarity with the institution, time, perceived legitimacy of seeking help, and expectations about whether support will be useful. Fong et al. \citep{fong2023} emphasize that help-seeking is not a unitary behaviour and that adaptive forms differ from avoidance or solution-seeking. Our administrative data cannot distinguish these mechanisms, so they should be investigated rather than assumed. Nevertheless, the observed mismatch between strain and formal help-seeking identifies a group---students with individual or compounded strain and little formal support use---for whom qualitative or survey research could be especially informative.

\subsection{Adaptation does not occur through a single channel}

The longitudinal results further show that support use and context selection are distinct adaptive channels. The simple hypothesis that students who do poorly in an assigned MU context later compensate by enrolling with higher-outcome Calculus instructors is not supported. If anything, individual-strain students subsequently appear with slightly lower historical-outcome instructor ranks. That result persists when the analysis is restricted to periods with multiple rankable instructors and stronger historical coverage.

The immediate post-failure transition is different. Nearly every student changes professor after a first MU failure, and most comparable students move toward professors associated with historically stronger prior outcomes. The pattern remains when historical context is measured by absolute prior pass rate or absolute prior mean grade, reducing concern that it is produced solely by changes in the composition of instructors across periods.

These two findings can coexist because they describe different decision points. The transition to Calculus occurs after students have already passed MU and progressed along the mathematics sequence. A failed MU attempt, by contrast, creates an immediate need to repeat the same course. Failure may change the salience of instructor context, the advice students receive, or the options presented during re-enrolment. The data cannot identify which mechanism is responsible, but the contrast suggests that academic adaptation is event-dependent rather than a stable tendency to seek an ``easier'' environment.

Importantly, movement in instructor context does not coincide with greater formal support use. CMAT attendance falls after the first failure, even though instructor switching is nearly universal. This decline cannot be interpreted as disengagement. The institutional setting includes changing incentive conditions around first-year support, and students may also substitute toward other resources or alter study strategies that are not observed. The result instead reinforces the conceptual value of keeping formal help-seeking and other adaptive behaviour separate.

\subsection{Repeated failure as a high-risk trajectory}

Students who fail repeatedly form a substantively important population rather than a nuisance group to be removed from analysis. The observed next-attempt pass rate declines from 64.6\% after the first failure to 54.5\% after the second and 40.0\% after the third. Later samples are small, but the trajectory illustrates accumulating academic risk.

At the same time, the formal models do not show that more severe numeric failure produces proportionally stronger adaptation. This is an important negative result. It suggests that simply increasing academic difficulty does not mechanically increase help-seeking or context change. For institutional practice, interventions targeted only on the assumption that students will respond adaptively once difficulty becomes sufficiently visible may therefore miss students who need support but do not mobilize it.

A stronger future design would combine administrative trajectories with measures of perceived difficulty, academic self-efficacy, help-seeking attitudes, and reasons for instructor enrolment. Such data could distinguish whether low support use under individual strain reflects avoidance, competing obligations, alternative sources of support, lack of awareness, or a belief that formal support is not appropriate. The current administrative design identifies the pattern but not the latent mechanism.

\subsection{Degree programme as disciplinary context}

Programme-level heterogeneity is statistically clear for first-MU CMAT use, later Calculus support use, and later instructor context, and the response to experience state varies across programmes. Yet the additional explained variance is modest. This balance is theoretically useful. Degree programme appears to be a meaningful structural context without becoming a deterministic student characteristic.

Kahu's sociocultural and structural framing provides a natural interpretation: disciplines can differ in mathematical demand, norms around help-seeking, sequencing of required mathematics, peer cultures, and the perceived value of mastering quantitative material \citep{kahu2013}. The present data cannot determine which of these mechanisms produces programme differences. Actuar\'ia, for example, displays high formal-support use without an unusually high average Calculus instructor-outcome rank, a pattern compatible with stronger institutionalized use of mathematics support but not proof of a distinct ``engagement culture.''

The programme results should therefore motivate targeted follow-up rather than programme stereotypes. Survey or qualitative work could test whether students in quantitatively intensive programmes perceive CMAT use as routine, whether peers and faculty normalize formal help-seeking differently, or whether curriculum demands simply create more opportunities to seek support.

